% Isolated 2026-09-15 discussion synchronization.
% Isolated 2026-09-14 writing revision; derived from the current checkout.
\textbf{What the evidence establishes.} Source diversity, executed
mechanisms, and one-run correctness differences are distinct measurements.
The Phase-I pattern did not recur in the revised Phase-II setting; multiple
protocol changes prevent identifying the cause. The corrected analysis of
the original gate shows no clear headroom improvement. Because that gate
also excludes allowed prompt strategies, it does not identify the effect
of verification alone.

\textbf{What remains unidentified.} Same-code reruns can disagree and
create oracle gains. Outcome flip rates measure instability, not the
standard error of a treatment contrast, and cannot define a numerical
noise floor for that contrast. A matched-cost clone control, independent
repeat validation, and validated calibration across mechanism types
(Appendix~\ref{app:gatecontrols})
are needed before attributing headroom to stable additional capability.


The post-review provider acquisition attempts did not yield a complete common pool: transport interruptions left completed responses, unclosed requests, and unstarted attempts in separate auditable states. These records are retained for provenance and missingness accounting, but are excluded from claims about population performance, cost, or stable complementarity. Thus, the manuscript does not use provider availability or partial acquisition as evidence for the unresolved clone-control question.

\textbf{Limitations.} The main comparison uses one text-to-SQL target and
six fixed builders, with three generation seeds per builder. Phase I is
discovery evidence with disclosed task overlap, a proxy judge, and a
post-trace human control. The new statistical analysis is post-review and
does not add independent executions. K-matched and R2 results now share the
corrected data and resampling chain. W1/W3 now share the canonical data
but remain descriptive; selection uncertainty and independent utility
validation are outstanding.
The one-seed MATH-500 free-arm probe provides parallel cross-domain
support: outcome diversity and one-run headroom reappear, but bare remains
the best fixed member. Its protocol is in Appendix~\ref{app:crossdomain};
it supplies no clone, repeat, repaired-gate, or stable-complementarity test.
